\chapter{Kết luận và hướng phát triển}

RBAC Guard kết nối hai lớp kiểm soát thường được tách rời: authorization trả
lời quyền tại thời điểm request, còn pipeline log đưa quyết định và hành vi vào
một đầu ra có thể truy vết. Demo Nova Bank làm rõ enforcement tại backend và
separation of duties; pipeline Python làm rõ parser, rule, scoring, context
finding và incident. Các artifact CSV/JSON cùng bộ test biến kịch bản demo
thành kết quả có thể tái lập thay vì chỉ là ảnh giao diện.

Giới hạn quan trọng là dữ liệu tổng hợp, regex SQL Injection cố định, ngưỡng
password guessing cố định theo user--IP và heuristic context đơn giản. Dataset
có nhãn được tạo để kiểm thử scenario mục tiêu, do đó metric đo mức độ phù hợp
với implementation hiện có chứ không phải chất lượng sản xuất. SQLite, session
token demo và policy baseline cũng chỉ phù hợp mục đích minh họa.

Các bước phát triển tiếp theo nên bắt đầu bằng log đã ẩn danh, có nhãn độc lập
và nhiều khoảng thời gian để đo false positive, false negative, độ trễ và chi
phí triage. Khi có bằng chứng vận hành, hệ thống có thể mở rộng schema/provenance
log, quản lý rule và ngưỡng theo cấu hình, streaming ingestion, tích hợp SIEM,
RBAC hierarchy hoặc mô hình phát hiện bổ sung. Mỗi bước cần giữ nguyên nguyên
tắc của nguyên mẫu: kết quả phải giải thích được và tái lập từ input công bố.
